% 猎户座（综述）
% license CCBYSA3
% type Wiki

本文根据 CC-BY-SA 协议转载翻译自维基百科\href{https://en.wikipedia.org/wiki/Orion_(constellation)}{相关文章}。

猎户座是在北天球冬季可见的一组显著恒星。它是现代划分的88个星座之一，同时也名列公元2世纪天文学家托勒密所列出的48个星座之中。其名称源自希腊神话中的一位猎人。

在北半球的冬季夜晚，猎户座最为醒目，另外还有五个星座与其一同构成“冬季六边形”星群。猎户座中最明亮的两颗恒星是参宿七（$\beta$）和参宿四（$\alpha$），二者均位列夜空中最亮的恒星之列；它们都是超巨星，并且亮度略有变化。此外，猎户座中还有六颗视星等亮于3.0的恒星，其中三颗构成了著名的猎户座腰带这一短而笔直的星群。猎户座还是一年一度猎户座流星雨的辐射点，该流星雨是与哈雷彗星相关的最强流星雨之一；同时，猎户座内还包含猎户座星云，这是夜空中最明亮的星云之一。
\subsection{特征}
\begin{figure}[ht]
\centering
\includegraphics[width=6cm]{./figures/c950af7d2bff8c26.png}
\caption{猎户座的星群示意图，如肉眼所见，并绘制了连线。} \label{fig_Orion_1}
\end{figure}
猎户座的西北方毗邻金牛座，西南方为波江座，南方为天兔座，东方为麒麟座，东北方为双子座。猎户座覆盖约594平方度，在88个星座中面积排名第26位。该星座的边界由比利时天文学家欧仁·德尔波特于1930年划定，其边界由一个26边形构成。在赤道坐标系中，这些边界的赤经范围介于 $04^{h}43.3^{m}$与$06^{h}25.5^{m}$之间，而赤纬范围则介于 22.87° 与 −10.97° 之间。$^{[2]}$ 国际天文学联合会于1922年采用的该星座三字母缩写为 “Ori”。$^{[3]}$

猎户座在一月至四月的夜空中最为清晰可见，$^{[4]}$ 此时北半球正值冬季，而南半球为夏季。在热带地区（距赤道约8°以内），猎户座会通过天顶。

从五月到七月（北半球的夏季、南半球的冬季），猎户座位于白昼天空中，因此在大多数纬度地区不可见。然而，在南半球冬季的数月中，南极洲大部分地区即使在正午太阳也位于地平线以下。此时，恒星（以及猎户座，但仅限于最亮的恒星）会在地方正午前后数小时内于暮光中可见，出现在北方低空、天空中最明亮的区域，太阳正位于地平线下方不远处。在同一时刻、南极点本身（阿蒙森—斯科特南极站），参宿七仅高出地平线8°，而猎户座腰带几乎沿着地平线掠过。南半球的夏季月份中，猎户座在其他地区通常可于夜空中观测，但在南极洲却完全不可见，因为在南极圈以南的这一时期太阳并不落下。$^{[5][6]}$

在接近赤道的国家（例如肯尼亚、印度尼西亚、哥伦比亚、厄瓜多尔），猎户座在12月午夜前后位于头顶上方，而在2月则出现在傍晚的天空中。
\subsection{导航用途}
\begin{figure}[ht]
\centering
\includegraphics[width=6cm]{./figures/96807e4588169a15.png}
\caption{利用猎户座定位邻近星座中的恒星} \label{fig_Orion_2}
\end{figure}
猎户座在定位其他恒星时极为有用。将猎户座腰带的连线向东南方向延伸，可以找到天狼星（$\alpha$ CMa）；向西北方向延伸，则可找到毕宿五（$\alpha$ Tau）。沿着猎户座两肩向东方画一条直线，可指示南河三（$\alpha$ CMi）的方向。从参宿七穿过参宿四所作的连线则指向双子座的北河二与北河三（$\alpha$ Gem 与 $\beta$ Gem）。此外，参宿七还是“冬季圆环”星群的一部分。天狼星和南河三也可以通过从猎户座出发的假想连线加以定位（见星图），它们同时也是“冬季三角形”和“冬季圆环”的组成恒星。$^{[7]}$
\subsection{特征}
\begin{figure}[ht]
\centering
\includegraphics[width=6cm]{./figures/fa98c4485118af07.png}
\caption{约1825年在伦敦出版的一套星图卡片《乌拉尼娅之镜》中所描绘的猎户座} \label{fig_Orion_3}
\end{figure}
猎户座中最明亮的七颗恒星在夜空中构成了一个极具辨识度的沙漏形星群或图案。其中四颗恒星——参宿七（Rigel）、参宿四（Betelgeuse）、参宿五（Bellatrix）和参宿六（Saiph）——形成一个近似矩形的大轮廓，而在其中心位置则是由参宿一（Alnitak）、参宿二（Alnilam）和参宿三（Mintaka）组成的猎户座腰带三星。猎户的头部由另一颗第八颗恒星标示，即觜宿一（Meissa），对观测者而言也相当明亮。自腰带向下延伸的是一条由三颗天体构成的较短直线，即猎户座之剑（其中间的“星体”实际上并非恒星，而是猎户座星云），亦称为猎人的佩剑。

其中许多恒星都是高光度、炽热的蓝色超巨星，腰带与佩剑中的恒星共同构成了猎户座 OB1 协会。参宿四以其显著的红色光芒格外醒目，但它可能仍是同一恒星集团中的一颗逃逸成员。
 
\subsubsection{亮星}
\begin{figure}[ht]
\centering
\includegraphics[width=6cm]{./figures/9cabe4a5ffcad42a.png}
\caption{按距离排列的猎户座恒星（红—绿三维视图）及各恒星的亮度（以恒星大小表示）} \label{fig_Orion_4}
\end{figure}
\begin{itemize}
\item 参宿四（Betelgeuse），又称猎户座$\alpha$星（Alpha Orionis），是一颗质量巨大的 M 型红超巨星，正接近其生命末期。它是猎户座中第二亮的恒星，同时也是一颗半规则变星。$^{[8]}$在传统描绘中，它位于猎人的右肩（假定猎人正面朝向观测者）。参宿四通常是夜空中第十一亮的恒星，但到2019年底，其亮度排名曾在第十至第二十三位之间变化。$^{[9][10]}$其生命终结预计将以一次超新星爆发告终，该事件从地球上将极为显著，甚至可能亮过月亮，并在白昼中可见。这一事件最有可能在未来十万年内发生。$^{[11]}$
\item 参宿七（Rigel），又称猎户座$\beta$星（Beta Orionis），是一颗 B 型蓝色超巨星，是夜空中第七亮的恒星。与参宿四类似，参宿七的核心正在进行重元素聚变，并将在天文学尺度上不久后结束其超巨星阶段，最终也可能以超新星爆发收场。它在星座形象中代表猎人的左脚。$^{[12]}$
\item 参宿五（Bellatrix）由约翰·拜耳标记为猎户座$\gamma$星（Gamma Orionis），是夜空中第27亮的恒星。参宿五被认为是一颗 B 型蓝色巨星，但其质量不足以发生超新星爆发。其高光度主要源于极高的表面温度，而非巨大的半径。$^{[13]}$参宿五标示猎户座的左肩，其名称意为“女战士”，有时也被俗称为“亚马逊星”。$^{[14]}$它是猎户座中距离太阳系最近的主要恒星，距离仅为244.6光年。
\item 参宿三（Mintaka）被指定为猎户座$\delta$星（Delta Orionis），尽管它是猎户座腰带三星中最暗的一颗。其名称意为“腰带”。它是一个多星系统，由一颗 B 型蓝色巨星和一颗质量更大的 O 型主序星组成。参宿三系统构成一对食变双星，其中一颗恒星掩食另一颗时会导致亮度下降。参宿三是猎户座腰带中最西的一颗，同时也是最北的一颗。$^{[7]}$
\item 参宿二（Alnilam）被指定为猎户座$\epsilon$星（Epsilon Orionis），其名称源自阿拉伯语，意为“珍珠串”。$^{[7]}$ 它是猎户座腰带三星中位于中央且最明亮的一颗。参宿二是一颗 B 型蓝色超巨星；尽管其距离太阳几乎是另外两颗腰带恒星的两倍，但其巨大光度使其视星等几乎相同。参宿二由于体积巨大而正在迅速失去质量。它是猎户座中距离最远的主要恒星，约为1344光年。
\item 参宿一（Alnitak），意为“束带”，被指定为猎户座$\zeta$星（Zeta Orionis），是猎户座腰带中最东的一颗。它是一个三合星系统，其中主星是一颗炽热的蓝色超巨星，也是夜空中最明亮的 O 型恒星。
\item 参宿六（Saiph）由拜耳标记为猎户座$\kappa$星（Kappa Orionis），在星座形象中代表猎户的右脚。它的距离和体积与参宿七相近，但看起来要暗得多。其名称意为“巨人的佩剑”。
\item 觜宿一（Meissa）被指定为猎户座$\lambda$星（Lambda Orionis），构成猎户座的头部，是一个多星系统，其合成视星等为3.33。其名称意为“闪耀者”。
\end{itemize}
\begin{figure}[ht]
\centering
\includegraphics[width=8cm]{./figures/5a047f4793781613.png}
\caption{猎户座最亮的8颗恒星} \label{fig_Orion_5}
\end{figure}
\subsubsection{腰带}  
猎户座腰带（Orion's Belt）是该星座中的一个星群。它由三颗明亮的恒星组成：参宿一（Alnitak，猎户座 $\zeta$ 星）、参宿二（Alnilam，猎户座 $\epsilon$ 星）和参宿三（Mintaka，猎户座 $\delta$ 星）。参宿一距离地球约800光年，光度约为太阳的10万倍，视星等为1.8；其辐射中有很大一部分位于人眼不可见的紫外波段。$^{[15]}$ 参宿二距离地球约2000光年，视星等为1.70，其紫外辐射光度约为太阳的37.5万倍。$^{[16]}$ 参宿三距离约915光年，视星等为2.21，其光度约为太阳的9万倍，并且是一对双星：两颗恒星以5.73天的周期相互绕转。$^{[17]}$ 在北半球，猎户座腰带在1月夜空中约晚上9点最为清晰可见，此时它大致位于当地子午线附近。$^{[18]}$

在参宿一的正西南方向不远处是 $\sigma$ 猎户座，这是一个由五颗恒星组成的多星系统，其合成视星等为3.7，距离约1150光年。在参宿三的西南方向则是四合星系统 $\eta$ 猎户座。
\subsubsection{佩剑}  
猎户座之剑包含猎户座星云、梅西耶43号星云、Sh~2--279（亦称“奔跑人星云”），以及 $\theta$ 猎户座、$\iota$ 猎户座和 42 猎户座等恒星。
\subsubsection{头部}  
三颗恒星构成一个小型三角形，用以标示猎户座的头部。其中顶点由觜宿一（Meissa，猎户座 $\lambda$ 星）标记，这是一颗光谱型为 O8~III 的炽热蓝色巨星，视星等为3.54，距离约1100光年。底边由 $\phi_1$ 猎户座与 $\phi_2$ 猎户座构成。其附近还有一颗年轻恒星 FU~猎户座。
\subsubsection{木棒}  
从参宿四向北延伸的一组恒星构成了猎户的木棒。$\mu$ 猎户座标示肘部，$\nu$ 与 $\xi$ 猎户座标示木棒的柄部，而 $\chi_1$ 与 $\chi_2$ 猎户座标示木棒的末端。在 $\chi_1$ 的正东方不远处，是一颗米拉型变星红巨星 U~猎户座。
\subsubsection{盾牌} 
在参宿五以西分布着六颗恒星，均被指定为 $\pi$ 猎户座（$\pi_1$ Ori、$\pi_2$ Ori、$\pi_3$ Ori、$\pi_4$ Ori、$\pi_5$ Ori 和 $\pi_6$ Ori），它们共同构成了猎户的盾牌。
\subsubsection{流星雨}  
每年约在10月20日前后，猎户座流星雨（Orionids）达到极大。该流星雨的辐射点位于与双子座交界的区域，每小时最多可观测到约20颗流星。其母体为哈雷彗星。$^{[19]}$
\subsubsection{深空天体}
\begin{figure}[ht]
\centering
\includegraphics[width=6cm]{./figures/4d6cf03ee91eb8e7.png}
\caption{这一视图突显了许多更为微弱的结构特征，例如巴纳德环。} \label{fig_Orion_6}
\end{figure}
\begin{figure}[ht]
\centering
\includegraphics[width=6cm]{./figures/05e5db3085b271fe.png}
\caption{北天窄带巡天所呈现的更宽视场（50° × 39°）图像，显示了猎户—波江超泡的大部分结构。} \label{fig_Orion_7}
\end{figure}
悬挂在猎户座腰带下方的是他的佩剑，由 $\theta_1$ 猎户座与 $\theta_2$ 猎户座这组多星系统组成，其中包含被称为“梯形星团”的恒星群以及猎户座星云（M42）。这是一个极为壮观的天体，即使用肉眼也能清楚辨认出它并非单一恒星。借助双筒望远镜，可以观测到其中孕育中的恒星、发光气体与尘埃云。梯形星团包含大量新生恒星，其中包括若干褐矮星，其距离大约为1500光年。该星团因四颗明亮恒星构成梯形而得名，其主要照明来源是这些仅有数十万年历史的年轻恒星。钱德拉 X 射线天文台的观测显示，主星的温度可高达 60000 开尔文，同时在周围星云中仍然存在活跃的恒星形成区域。$^{[20]}$

M78（NGC 2068）是位于猎户座内的一座星云。其整体视星等为8.0，明显暗于其南侧的大猎户座星云；然而，它们距离地球大致相同，约为1600光年。在望远镜目镜中，M78 很容易被误认为是一颗彗星。M78 与变星 V351 猎户座相关联，该恒星的亮度变化可在极短时间尺度内被观测到。$^{[21]}$猎户座内另一座较为明亮的星云是 NGC 1999，同样靠近大猎户座星云，其综合视星等为10.5，距离地球约1500光年。变星 V380 猎户座嵌置于 NGC 1999 之中。$^{[22]}$

另一座著名的星云是位于参宿一（猎户座$\zeta$星）附近的 IC 434，即马头星云。该星云包含一团暗尘埃云，其形状正是该星云名称的由来。

NGC 2174 是一座发射星云，距离地球约6400光年。

除上述星云之外，使用小型望远镜对猎户座进行巡察，还可以发现大量引人入胜的深空天体，包括 M43、M78，以及多星系统如$\iota$猎户座和$\sigma$猎户座。使用更大口径的望远镜，则可能观测到火焰星云（NGC 2024），以及更暗、更紧密的多星系统和星云。在极其黑暗的夜晚或借助长时间曝光摄影，还可以看到巴纳德环。

上述所有星云均属于更大的猎户座分子云复合体，该复合体位于约1500光年之外，尺度达数百光年。由于其距离较近，它是从地球可见的最为活跃的恒星形成区域之一。猎户座分子云复合体构成了一个更为宏大的结构——猎户—波江超泡的东部区域，该结构可在 X 射线和氢发射中观测到。
\subsection{历史与神话} 
猎户座极具辨识度的形态在世界各地的众多文化中均被认出，并与大量神话传说相联系。在现代世界中，猎户座也被用作象征符号。
\subsubsection{古代近东}
\begin{figure}[ht]
\centering
\includegraphics[width=6cm]{./figures/9778110b3b9e0b4e.png}
\caption{猎户座的艺术描绘} \label{fig_Orion_8}
\end{figure}
青铜时代晚期的巴比伦星表将猎户座称为 MULSIPA.ZI.AN.NA$^{[note\ 1]}$，意为“天上的牧羊人”或“阿努的真正牧羊人”——其中阿努是天界的主神。$^{[23]}$在巴比伦体系中，该星座被视为神圣之物，归属于帕普舒卡尔（Papshukal）和宁舒布尔（Ninshubur），二者皆为履行“诸神信使”职能的次要神祇。帕普舒卡尔与巴比伦界碑上“行走之鸟”的形象密切相关，而在星图中，“公鸡”的形象位于“真正的牧羊人”形象的下方和后方——这两个星座分别以鸟形与人形象征诸神的报信者。$^{[24]}$

在古埃及，猎户座中的恒星被视为一位名为 Sah 的神祇。由于猎户座升起早于天狼星，而天狼星的偕日升是埃及太阳历的基础，因此 Sah 与人格化天狼星的女神 Sopdet 密切相关。神祇 Sopdu 被认为是 Sah 与 Sopdet 之子。Sah 后来与奥西里斯（Osiris）发生神格融合，而 Sopdet 则与奥西里斯的神话妻子伊西斯（Isis）融合。在公元前24至23世纪的《金字塔文》中，Sah 是诸多神祇之一，亡故的法老被认为会在来世化身为这些神祇的形态。$^{[25]}$

亚美尼亚人将其传说中的族长与建国者海克（Hayk）与猎户座等同。Hayk 亦是《圣经》亚美尼亚文译本中猎户座的名称。$^{[26]}$

《圣经》中三次提及猎户座，将其称为 “Kesil”（כסיל，字面意为“愚者”）。然而，该名称在词源学上可能与“基斯流月”（Kislev）有关，即希伯来历的第九个月（约为公历11月至12月），而“基斯流”又可能源自希伯来语词根 K-S-L，如 “kesel, kisla”（כֵּסֶל, כִּסְלָה，意为“希望、积极性”），即对冬季降雨的期盼。相关经文包括：《约伯记》9:9（“他造了北斗、猎户”）、38:31（“你能解开猎户的带吗？”）以及《阿摩司书》5:8（“造昴星和猎户的那一位”）。

在古代亚兰文化中，该星座被称为 Nephîlā′，而尼非林（Nephilim）则被认为是猎户的后裔。$^{[27]}$
\subsubsection{希腊—罗马古代}  
在希腊神话中，猎户（Orion）是一位体型巨大、具有超自然力量的猎人，$^{[28]}$出生于蛇发女妖欧律阿勒（Euryale）与海神波塞冬（罗马名尼普顿）。有一种神话讲述了大地女神盖娅因猎户狂妄宣称要杀尽世上一切动物而大怒，遂派遣天蝎去惩罚他。这被视为天蝎座与猎户座从不在同一时间出现在天空中的原因。然而，持蛇者俄斐俄古斯（Ophiuchus）以解毒药复活了猎户，因此其星座被认为位于天蝎与猎人之间的天区。$^{[29]}$

猎户座亦见于多部古典文学作品中，包括贺拉斯的《颂歌》（第3卷第27首第18行）、荷马的《奥德赛》（第5卷第283行）与《伊利亚特》，以及维吉尔的《埃涅阿斯纪》（第1卷第535行）。
\subsubsection{中东}
\begin{figure}[ht]
\centering
\includegraphics[width=6cm]{./figures/2e6873de2c971d01.png}
\caption{如公元962年的波斯天文学著作《定星书》中所描绘的那样。在这一表现形式中，猎户座被绘制在天球仪上，因此与其在真实天空中的外观相比呈现出左右相反的效果。} \label{fig_Orion_9}
\end{figure}
在中世纪伊斯兰天文学中，猎户座被称为 al-jabbar，意为“巨人”。$^{[30]}$ 猎户座中第六亮的恒星参宿六（Saiph），其名称源自阿拉伯语 saif al-jabbar，意为“巨人的佩剑”。$^{[31]}$
\subsubsection{中国}  
在中国，猎户座是二十八宿之一，称为参宿（Sieu / Xiù）。其名为“参”（Shen），字面意为“三”，指的是猎户座腰带的三星。汉字“參”（拼音：shēn）最初即指猎户座（中文：參宿；拼音：shēnxiù）；其距今三千多年的商代字形，在字的上部描绘了三颗位于人头之上的恒星，象征猎户座腰带（下部表示读音的构件为后期所加）。$^{[32]}$

\subsubsection{印度}  
《梨俱吠陀》中将该星座称为 Mriga（鹿）。$^{[33]}$

“舞王湿婆”（Nataraja，“宇宙之舞者”）常被解读为猎户座的象征。梨俱吠陀中的鲁德罗（Rudra，即湿婆的早期形态）是印度占星术中阿尔德拉宿（Ardra nakshatra，对应参宿四 Betelgeuse）的主宰神。$^{[34]}$

印度乌代吉里与坎达吉里洞穴中、公元前1世纪雕刻的耆那教符号$^{[35]}$，在外形上与猎户座呈现出显著相似性。

布吉斯（Bugis）水手将猎户座腰带的三颗星称为 tanra tellué，意为“三之标志”。$^{[36]}$
\subsubsection{欧洲民间传统}  
在古老的匈牙利传统中，猎户座被称为“弓箭手”（Íjász）或“收割者”（Kaszás）。在近年重新发现的神话中，他被称为宁录（Nimrod，匈牙利语：Nimród），是最伟大的猎人，也是胡诺尔（Hunor）与马戈尔（Magor）这对双胞胎的父亲。右上方的 $\pi$ 星与 $o$ 星共同构成反曲的弓或高举的镰刀。在其他匈牙利传统中，猎户座腰带被称为“法官之杖”（Bírópálca）。$^{[37]}$

在爱尔兰与苏格兰，猎户座被称为 An Bodach，这是爱尔兰民间传说中的一个形象，其名称字面意为“有阳具之人”，被视为女巫 Cailleach 的丈夫。$^{[38][39][40]}$

在斯堪的纳维亚传统中，猎户座腰带被称为“弗丽嘉的纺锤”（Frigg's Distaff，friggerock）或“芙蕾雅的纺锤”。$^{[41]}$

芬兰人将猎户座腰带及其下方的恒星称为“维南莫宁的镰刀”（Väinämöisen viikate）。$^{[42]}$ 参宿二、参宿一与参宿三组成的星群还被称为“维南莫宁的腰带”（Väinämöisen vyö），而从腰带“垂下”的恒星则被称为“卡勒瓦之剑”（Kalevanmiekka）。

在西伯利亚，楚科奇人将猎户座视为一名猎人；他射出的箭由毕宿五（Aldebaran，金牛座 $\alpha$ 星）所代表，其整体形象与其他西方传统中的描绘相同。$^{[43]}$

大众媒体中曾有人声称，德国盖森克勒斯特洞穴中一件距今约35000至40000年的象牙雕刻“朝拜者”（Adorant）是已知最早的猎户座形象描绘。然而学界否定了这一解读，指出所谓的“腰带”与“佩剑”等细节，很可能源自象牙材料本身的纹理结构。$^{[44][45][46][47]}$
\subsubsection{美洲}  
墨西哥西北部的塞里人（Seri）将猎户座腰带的三颗星称为 Hapj（意为“猎人”），由三颗星组成：Hap（骡鹿）、Haamoja（叉角羚）与 Mojet（大角羊）。其中 Hap 位于中央，被猎人射中，其血滴落在蒂布龙岛上。$^{[48]}$

在西班牙及大多数拉丁美洲地区，这三颗星被称为“Las tres Marías”（“三位玛利亚”）。在波多黎各，它们被称为“Los Tres Reyes Magos”（“三位东方贤士”）。$^{[49]}$

奥吉布瓦／奇佩瓦原住民称该星座为 Mesabi，意为“伟大之人”。$^{[50]}$

在拉科塔原住民的传统中，Tayamnicankhu（猎户座腰带）被视为一头野牛的脊柱。猎户座的大矩形代表野牛的肋骨；邻近金牛座的昴宿星团是野牛的头部；而位于大犬座的天狼星（称为 Tayamnisinte）则是其尾巴。另一个拉科塔神话认为，猎户座下半部，即“手之星座”，代表一位酋长的手臂，该手臂因其自私而被雷霆之民撕裂，作为诸神的惩罚。其女儿承诺嫁给能够从天空中取回手臂之人，于是年轻的战士“坠星”（Fallen Star，其父为恒星、母为凡人）取回了手臂并迎娶了她，象征着在年轻一代的帮助下，诸神与人类之间重新达成和谐。食指由参宿七（Rigel）表示；猎户座星云代表拇指；猎户座腰带代表手腕；而波江座 $\beta$ 星则代表小指。$^{[51]}$
\subsubsection{南岛语族}  
猎户座的七颗主要恒星构成了波利尼西亚星座 Heiheionakeiki，该星座象征一种儿童用的绳结图形，类似于“翻花绳”。$^{[52]}$ 一些前殖民时期的菲律宾族群尤其将猎户座腰带区域称为“balatik”（弩机），因为其形状类似同名的陷阱，这种装置可自行发射箭矢，通常用于在灌木丛中捕捉野猪。$^{[53][54][55][56]}$ 后来的西班牙殖民统治使部分族群将猎户座腰带称为 “Tres Marias” 或 “Tatlong Maria”。$^{[57]}$

在毛利传统中，参宿七（Rigel，称为 Puanga 或 Puaka）与 Matariki 的庆典密切相关。隆冬时节，Matariki（昴宿星团）与参宿七在日出前的升起，标志着毛利新年的开始。$^{[58]}$

在爪哇文化中，该星座常被称为 Lintang Waluku 或 Bintang Bajak，意指水田犁具的形状。$^{[59]}$
\subsubsection{当代象征}  
腰带与佩剑的意象进入了现代西方流行文化，例如在两次世界大战期间，美国陆军第27步兵师的肩章徽记中即采用了这一形象，这可能源于该师首任指挥官约翰·F·奥莱恩少将（John F. O'Ryan）姓名的谐音双关。$^{[60][61]}$

电影发行公司猎户座影业（Orion Pictures）亦以该星座作为其标志。$^{[62]}$
\subsubsection{描绘}
\begin{figure}[ht]
\centering
\includegraphics[width=6cm]{./figures/35478d69812c7ca6.png}
\caption{9世纪《莱顿阿拉托斯》中所描绘的猎户座} \label{fig_Orion_10}
\end{figure}
在艺术描绘中，周围的星座有时会与猎户座产生关联：他常被描绘为站立在波江座旁，身边伴随着他的两只猎犬——大犬座与小犬座，并与金牛座搏斗。有时他被描绘成正在猎捕天兔座；也有时他手中拿着一张狮皮。

猎户座还有其他不同的视觉化方式。从南半球观察时，猎户座呈现为“南方在上”的朝向，而其腰带与佩剑在澳大利亚和新西兰常被称为“平底锅”或“汤锅”。在南非，讲南非荷兰语的人将猎户座腰带称为 Drie Konings（“三王”）或 Drie Susters（“三姐妹”）$^{[63]}$，而在都德的《磨坊来信》（Lettres de Mon Moulin，1866）中则被称为 les Trois Rois（“三王”）。Driekoningen（“三王”）这一称呼也常见于17至18世纪的荷兰星图与航海指南中。同样的三颗恒星在西班牙、拉丁美洲与菲律宾被称为 “Las Tres Marías”（“三位玛利亚”），而在波多黎各则称为 “Los Tres Reyes Magos”（“三位东方贤士”）。$^{[49]}$

即便是在传统描绘中，猎户座的形象也存在很大差异。西塞罗所绘制的猎户形象与现代描绘较为相近：猎人右手高举一张不明动物的皮，其手由猎户座$\omicron^2$星表示，而兽皮则由五颗 $\pi$猎户座恒星表示。参宿六（Saiph）与参宿七（Rigel）分别代表他的左膝与右膝，而 $\eta$猎户座与$\lambda$天兔座则分别代表他的左脚与右脚。如同现代描绘，参宿三（Mintaka）、参宿二（Alnilam）与参宿一（Alnitak）构成他的腰带。他的左肩由参宿四（Betelgeuse）表示，而$\mu$猎户座构成其左臂；觜宿一（Meissa）是他的头部，参宿五（Bellatrix）则是他的右肩。

希吉努斯（Hyginus）的描绘与西塞罗相似，但在若干关键细节上有所不同。西塞罗所绘的兽皮在希吉努斯的版本中变成了盾牌（由$\omicron$与$\pi$猎户座表示），而左臂不再由$\mu$猎户座标示，而是改为手持一根木棒（由$\chi$猎户座表示）。他的右腿由$\theta$猎户座表示，而左腿则由$\lambda$、$\mu$与$\epsilon$天兔座恒星共同构成。此后，西欧与阿拉伯世界的猎户座描绘大多沿袭了这两种模式。$^{[43]}$
\subsection{未来}
\begin{figure}[ht]
\centering
\includegraphics[width=6cm]{./figures/e0e94b9e231d8272.png}
\caption{动画展示了猎户座从公元前50000年至公元50000年的自行运动变化，其中 $\pi^3$ 猎户座的运动速度最快。} \label{fig_Orion_11}
\end{figure}
猎户座位于天球赤道之上，但由于地轴进动的影响，它并非始终处于这一位置。猎户座整体位于黄道以南，而它之所以恰好落在天球赤道附近，是因为与夏至对应的黄道点接近双子座与金牛座的交界处，位于猎户座以北。随着岁差的持续作用，猎户座最终将逐渐向南移动；到公元14000年时，猎户座将南移到足以使其从大不列颠纬度范围内完全不可见。$^{[64]}$

在更遥远的未来，由于自行运动的影响，猎户座中的恒星将逐渐偏离当前的星座形态。然而，从天文学尺度来看，猎户座中最明亮的恒星距离地球都非常遥远——例如远远超过天狼星的距离。因此，在由相对较近恒星构成的其他大多数星座因自行运动而逐渐扭曲成全新形态之后，猎户座仍将长期保持可辨识性。唯一的例外是其中部分恒星最终会以超新星形式爆发，例如参宿四（Betelgeuse），其预计将在未来一百万年内发生爆炸。$^{[65]}$

\subsection{参见}  
\begin{itemize}
\item EURion 星座  
\item 《哈勃3D》（2010），一部包含对猎户座星云进行精细 CGI “飞行穿越”的 IMAX 电影  
\item 猎户相关理论  
\item 猎户（神话）  
\item 猎户（中国天文学）  
\item 猎户复合体的外围星云区域  
\item Aurvandill  
\item Glooscap  
\item Heiheionakeiki  
\item Julpan  
\item Nataraja  
\item Osiris  
\item Papsukkal  
\item Urania  
\item 冬季六边形  
\end{itemize}
\subsection{参考文献}  
\subsubsection{说明性注释}  
\begin{enumerate}
\item 在这些列表中，用于表示“星座”或“恒星”的限定符号为 MUL（𒀯）。参见巴比伦星表。
\end{enumerate}
\subsubsection{引用}
\begin{enumerate}
\item 柯克帕特里克，J·戴维；马罗科，费德里科等（2024年4月）。“基于全天空20秒差距范围内约3600颗恒星和棕矮星普查的初始质量函数”. 天体物理学报增刊系列. 271 (2): 55. arXiv:2312.03639. Bibcode:2024ApJS..271...55K. doi:10.3847/1538-4365/ad24e2.
\item “猎户座，星座边界”。星座。国际天文学联合会。于2013年6月4日从原文存档。检索日期：2013年3月22日
\item 拉塞尔，亨利·诺里斯（1922）。“星座的新国际符号”。大众天文学。第30卷。第469–71. 文献标识码:1922PA.....30..469R.
\item 埃利亚德，大卫；蒂里昂，威尔（2008）[1993]。《南天指南》（第3版）。维多利亚州波特墨尔本：剑桥大学出版社。第4页。ISBN 978-0-521-71405-1.
\item “南半球天空入门指南”. dibonsmith.com.
\item “南半球晚间星空图”. skymaps.com.
\item 斯塔尔 1988, 第61页。
\item 本月变星，猎户座α星.本季变星美国变星观测者协会。2000年。于2009年1月22日从原文存档。检索日期：2009-02-26
\item 迪金森，大卫（2019-12-26）。“等待参宿四：这颗暴躁的恒星怎么了？”. 今日宇宙。检索日期：2025-05-04.
\item 多兰，克里斯。“参宿四”。于2011年11月24日存档自原文。检索日期：2023年8月6日。
\item 普赖尔，瑞安（2019年12月26日）.“一颗巨大的红巨星行为异常，科学家认为它可能即将爆炸”. CNN.
\item “参宿七”. 吉姆·凯勒的恒星. 伊利诺伊大学厄巴纳-香槟分校。2009年。于2009年2月22日从原文存档。检索日期：2009-02-26.
\item “贝拉特里克斯”。吉姆·凯勒的恒星伊利诺伊大学厄巴纳-香槟分校。2009年。于2009年2月22日从原文存档。检索日期：2009-02-26
\item 多兰，克里斯。“贝拉特里克斯”。于2011年11月30日存档自原文。检索日期：2023年8月6日。
\item "Alnitak". Stars.astro.illinois.edu. Retrieved 2012-05-16.
\item “阿尔尼拉姆”。吉姆·凯勒的恒星伊利诺伊大学厄巴纳-香槟分校2009年于2011年11月24日存档自原始网页检索日期：2011年11月28日
\item “明塔卡”。吉姆·凯勒的恒星伊利诺伊大学厄巴纳-香槟分校2009年。于2011年11月24日从原始网页存档。检索日期：2011年11月28日
 Dolan, Chris. "Orion". Archived from the original on 2011-12-07. Retrieved 2011-11-28.
 詹尼斯肯斯，彼得（2012年9月）。“绘制流星体轨道：新流星雨被发现”。天空与望远镜。第22页。
 威尔金斯，杰米；邓恩，罗伯特（2006）。300个天文天体：宇宙视觉参考（第1版）。纽约布法罗：萤火虫图书。ISBN 978-1-55407-175-3.
 列维 2005, 第99–100页。
 莱维 2005, 第107页。
 罗杰斯，约翰·H.（1998）。“古代星座的起源：I.美索不达米亚传统”。英国天文协会期刊108：9–28。文献标识码1998JBAA..108....9R
 巴比伦星象学加文·怀特著，索拉里亚出版社，2008年，第218页及后续内容与第170页
 威尔金森，理查德·H.（2003）。古埃及诸神与女神全书。泰晤士与哈德逊出版社。第127页、第211页
 库尔基扬，瓦汉 (1968). “亚美尼亚史”. uchicago.edu. 密歇根州. 第8页
 皮克对圣经的注释, 1962年，第260页，第221f节。
 “星语——猎户座”. www.ianridpath.com.
 斯塔尔 1988, 第61–62页。
 梅特利茨基，多萝西 (1977). 中世纪英格兰的阿拉比之物。美国：耶鲁大学出版社。第79页。ISBN 0-300-11410-9.
 卡勒，詹姆斯·B.，“参宿四（猎户座κ）”, 恒星, 伊利诺伊大学, 检索日期2012-01-27
 汉语大字典（中文版），1992年（第163页）。湖北辞书出版社和四川辞书出版社出版，由台湾台北建宏出版社以繁体字形式再版；ISBN 957-813-478-9
 霍莱，P. V.（1998）。“吠陀天文学家”。印度天文学会会刊。26：91–106。文献标识码1998BASI...26...91HDOI10.1080/1468936042000282726821S2CID26503807
 斯里尼瓦桑，莎拉达（1998）。“吠陀天文学家”。世界考古学. 36: 432–50. 文献标识码:1998BASI...26...91H. DOI:10.1080/1468936042000282726821. S2CID 26503807.
 “必看”的印度遗产.asimustsee.nic.in。检索日期：2020-04-17
 凯利，大卫·H.；米隆，尤金·F.；阿维尼，A.F.（2011）。探索古代星空：古代与文化天文学概览。纽约，纽约：施普林格出版社。第344页。ISBN 978-1-4419-7623-9.
 托罗茨凯-维甘德·埃德：《老星》（“Old stars”），匈牙利（1915年）由技术图书出版社METRUM再版（1988年）。
 “一些爱尔兰语中的太空词汇 | Cúpla Focail Spáis i nGaeilge”，2025年3月13日。
 "The Bodach is a creature from Scottish folklore married to the Cailleach". September 21, 2021.
 Moffet, Amanda (April 19, 2011). "The Scottish Bodach -".
 舍恩，埃贝。（2004）。阿萨-托尔的锤子：信仰与传统中的神祇与巨人。费尔特与海斯勒，韦尔纳莫。第228页。
 埃洛，伊斯莫。“恒星与星座”。Ursa.fi。检索日期：2013年10月16日。
 斯塔尔 1988, 第63页。
 牛津史前雕像手册, 蒂莫西·英索尔 编，2017年，牛津大学出版社,ISBN 0199675619, 9780199675616, 谷歌图书
 拉彭格吕克，迈克尔（2001）。“天空中的类人生物：一块距今3.2万年的象牙盘是否展示了猎户座星座与妊娠日历的结合”。符号、日历与方位：文化中的天文遗产。欧洲文化天文学学会（SEAC）第九届年会。乌普萨拉天文台. 第51–55页.
 “德国盖森克勒特勒的装饰盘”.联合国教科文组织：天文遗产之门. 检索日期：2014年2月26日
 Whitehouse, David (January 21, 2003). "'Oldest star chart' found". BBC. Retrieved 26 February 2014.
 莫瑟，玛丽·B.；马雷特，斯蒂芬·A.（2005）。科姆卡克语词典：塞里语—西班牙语—英语 （PDF）（西班牙语和英语版）。索诺拉州埃尔莫西略市与墨西哥城：索诺拉大学与普拉萨-瓦尔德斯出版社。
 “首页 – 新日报”. Elnuevodia.com。于2013年10月24日从原文存档。检索日期：2013年10月16日。”
 Discover the Range, Explore Minnesota, 2025 [1]
 “通往宇宙的窗口”. Windows2universe.org. 检索日期：2017年1月13日.
 “夏威夷星际航线”.
 斯科特，威廉·亨利（1994）。巴兰盖：16世纪菲律宾的文化与社会。菲律宾马尼拉奎松市：马尼拉雅典耀大学出版社。第124页。ISBN 9789715501354.
 恩卡纳西翁，胡安·费利克斯（1885）。比萨亚语-西班牙语词典[印刷文本]（西班牙语和宿雾语）。第30页。
 “BALATIK：菲律宾人的本土之星 | 菲律宾社会科学评论”。于2023-07-27从原文存档。检索日期：2023-07-27。”
 “菲律宾民族天文学：群星如何指引我们的祖先——FlipScience”2020年12月18日
 “巴拉蒂克”.
 “普安加：预示马塔里基的星辰”.高等教育工会/特豪图卡胡兰吉. 于2022年1月27日从原文存档. 于2022年5月10日检索
 "Mengenal lintang waluku, rasi bintang fenomenal yang menjadi pranata mangsa". DetikJogja. Retrieved 16 July 2024.
 哈特，阿尔伯特·布什内尔（1920）。哈珀世界大战图画图书馆，第5卷。纽约：哈珀兄弟出版社。第358页。
 莫斯，詹姆斯·阿尔弗雷德；霍兰德，哈里·塞缪尔（1920）。《美国的战役：附法国与比利时美军战场指南。威斯康星州梅纳沙：乔治·班塔出版公司，第555页。
 金宇赫（2012-09-21）。“从山到月：10个电影制片厂标志及其背后的故事”. Time.com。检索日期：2015-09-22.
 “三王与开普云：两个天文谜题”.psychohistorian.org。于2010-01-29从原文存档。检索日期：2009-06-27
 “岁差”。Myweb.tiscali.co.uk。于2018-07-23从原始网页存档。检索日期：2012-05-16。
 威尔金斯，阿拉斯代尔（2011年1月20日）。“地球可能很快将迎来第二颗太阳”。io9。太空色情。于2015年11月14日存档自原文。检索日期：2012年2月5日
\end{enumerate}
\subsubsection{参考目录}
\begin{itemize}
\item 莱维，大卫·H.（2005）。深空天体。普罗米修斯图书。ISBN 1-59102-361-0.
\item 伊恩·里德帕斯和威尔·蒂里昂（2007）。恒星与行星指南, 科林斯，英国伦敦。ISBN 978-0-00-725120-9。普林斯顿大学出版社，新泽西州普林斯顿。ISBN 978-0-691-13556-4.
\item 斯塔尔，尤利乌斯·D·W.（1988）天空中的新图案, 麦当劳与伍德沃德出版公司,ISBN 0-939923-04-1.
\end{itemize}